\documentclass[11pt]{article}

\usepackage[utf8]{inputenc}
\usepackage[T1]{fontenc}
\usepackage[ngerman]{babel}
\usepackage[a4paper,margin=0.85in]{geometry}
\usepackage{graphicx}
\usepackage{array}
\usepackage{parskip}
\usepackage{ragged2e}

\pagestyle{empty}
\renewcommand{\baselinestretch}{1.0}

\begin{document}

% ---------- Title ----------
\begin{center}
    {\Large \textbf{Bewerbung}} \\[3cm]
    
    {\Huge \textbf{Zhongli Wang}} \\[0.4cm]
    
    {\huge Bewerbung als Data Scientist}
\end{center}

\vspace{0.6cm}

% ---------- Large Picture in the Middle ----------
\begin{center}
    \includegraphics[
        width=0.55\textwidth,
        height=0.45\textheight,
        keepaspectratio
    ]{../zhongli-wang-profile.jpg}
\end{center}

\vspace{0.7cm}

% ---------- Personal Details ----------
\begin{center}
\renewcommand{\arraystretch}{1.35}
\begin{tabular}{>{\bfseries}p{4.2cm} p{8.5cm}}

Anschrift: & Astastrasse 22, 33617 Bielefeld, Deutschland \\

Telefon: & +49 160 9541 6015 \\

E-Mail: & nkwangzhongli@gmail.com \\

\end{tabular}
\end{center}

\vfill

\newpage

\begin{flushright}
Zhongli Wang \\
Astastrasse 22 \\
33617 Bielefeld \\
Deutschland \\
\end{flushright}

\vspace{0.1cm}

\noindent
Recruiting-Team \\
Würth \\
Künzelsau, Deutschland \\

\vspace{0.05cm}

\noindent
\textbf{Bewerbung als Data Scientist} 

\noindent Sehr geehrtes Recruiting-Team,

ich schlie\ss{}e derzeit meine Promotion in Wirtschaftswissenschaft an der Universität Bielefeld ab (Abschluss September 2026). Meine Forschung verbindet quantitative ökonomische Modellierung, Datenanalyse, Simulationen und empirische Methoden. Die Position als Data Scientist bei Würth interessiert mich sehr, da sie mir die Möglichkeit bietet, datengetriebene Lösungen zu entwickeln, die Fachbereiche in ihrer täglichen Arbeit konkret unterstützen und echten Mehrwert schaffen.

Mein Hintergrund qualifiziert mich besonders für diese Rolle: Ich übersetze Businessfragen eigenständig in analytische und Modellierungsaufgaben und wähle dabei bewusst den einfachsten tragfähigen Lösungsweg. Im Rahmen meiner Promotion habe ich agentenbasierte Modelle zur Analyse von Wettbewerbsdynamiken und Investitionsentscheidungen entwickelt und validiert. Ich habe Python-basierte Prototypen und Produktionssysteme gebaut, die Datenverarbeitung, maschinelles Lernen (XGBoost, Random Forest, Neural Networks) und statistische Analyse kombinieren. Gleichzeitig habe ich komplexe ökonomische Konzepte in klare, entscheidungsrelevante Erkenntnisse für unterschiedliche Zielgruppen übersetzt – eine Fähigkeit, die für die Kommunikation zwischen Arbeitsebene und Management essentiell ist.

Meine technischen Fähigkeiten umfassen Python (pandas, NumPy, scikit-learn), SQL, statistische Modellierung und maschinelles Lernen. Ich habe Erfahrung mit LLMs, Prompt Engineering und agentenbasierten Workflows. Ich bin neugierig, denke nutzerorientiert und bin motiviert, Datenprodukte nicht nur zu entwickeln, sondern auch bis in die Anwendung zu bringen. Für meine Themen übernehme ich Verantwortung, erkenne Blocker frühzeitig und bringe aktiv Lösungsvorschläge ein.

Würth ist für mich besonders attraktiv, weil ich in einem Umfeld arbeiten möchte, in dem Datenprodukte echten Nutzen stiften und nicht in der Schublade landen. Ich bin motiviert, mit Fachbereichen zusammenzuarbeiten, ihre Anforderungen zu verstehen und analytische Werkzeuge zu entwickeln, die ihre Arbeitsweise spürbar unterstützen. Die Möglichkeit, an der Schnittstelle von Data Science, Produktdenken und Business zukunftsorientierte Lösungen zu gestalten, reizt mich sehr.

Ich freue mich sehr, in einem persönlichen Gespräch zu erläutern, wie meine Fähigkeiten zur Entwicklung von Datenprodukte bei Würth beitragen können.

\vspace{0.15cm}

\noindent Mit freundlichen Grü\ss{}en, \\
\vspace{2mm}\\
\noindent Zhongli Wang

\vspace{0.05cm}
\noindent\includegraphics[width=2.7cm,keepaspectratio]{../Wang-sig.jpg}

\end{document}
